\documentclass{article}

\usepackage{microtype}
\usepackage[accepted]{icml2025}

\usepackage{graphicx}
\usepackage{amsmath}
\usepackage{amssymb}
\usepackage{booktabs}

\icmltitlerunning{Do Tabular Foundation Models Scale with Dimensionality?}

\begin{document}

\twocolumn[
\icmltitle{Do Tabular Foundation Models Scale with Dimensionality? \\
           A Synthetic-Data Study from 10 to 10,000 Features}

\begin{icmlauthorlist}
\icmlauthor{Timothy Daley}{indep}
\end{icmlauthorlist}

\icmlaffiliation{indep}{Independent research}
\icmlcorrespondingauthor{Timothy Daley}{timy.pdaley@gmail.com}

\icmlkeywords{tabular data, foundation models, in-context learning, TabPFN, dimensionality, regression}

\vskip 0.3in
]

\printAffiliationsAndNotice{}

\begin{abstract}
TODO: fill in after results are final. One paragraph: motivation, method (synthetic
low-dimensional nonlinear signal embedded in up to 10,000 features, feature-bagging
ensemble for models beyond their native feature cap), and the headline empirical
finding (how TabPFN v2 / TabICL test R\textsuperscript{2} degrades relative to tree
baselines as D grows, under fixed-N and N-scaling-with-D regimes).
\end{abstract}

\section{Introduction}
\label{sec:intro}

Tabular foundation models -- transformer architectures pretrained on large corpora of
synthetic or real tabular datasets and applied zero-shot via in-context learning --
have shown striking accuracy on small- to medium-sized tabular regression and
classification benchmarks \citep{muller2022transformers, hollmann2025tabpfn}. Unlike
gradient-boosted trees or linear models, these models are not refit per dataset;
instead, the training set is passed as context alongside the query points, and a
single forward pass produces predictions.

This convenience comes with an architectural constraint that has received little
systematic study: every released tabular foundation model has a hard cap on the
number of features and rows it was pretrained and validated on. TabPFN v2
\citep{hollmann2025tabpfn} was trained and evaluated up to 500 features and 10,000
rows; TabICL \citep{qu2025tabicl} extends this considerably but still caps out in the
low thousands of features. Real tabular data routinely exceeds these limits --
genomics, click-stream, and sensor-fusion datasets can have thousands to tens of
thousands of columns -- so understanding \emph{how} these models degrade (or don't)
past their trained regime, and whether standard workarounds close the gap, is a
practical question for anyone considering a tabular foundation model outside a
benchmark setting.

We study this question with a controlled synthetic data generating process: a fixed,
low-dimensional (5-feature) nonlinear function determines the regression target, and
the remaining $D-5$ ambient features are either noisy nonlinear functions of the true
features or pure noise, with $D$ swept on a log-spaced grid from 10 to 10,000. Because
the true signal's dimensionality is fixed while the ambient dimensionality grows, this
isolates the effect of dimensionality \emph{per se} from the effect of adding more true
signal, which is the confound present in most real high-dimensional tabular benchmarks.
We compare TabPFN v2 and TabICL, each wrapped in a feature-bagging ensemble
\citep{feuer2023scaling} once $D$ exceeds their native cap, against Ridge regression,
Random Forest \citep{breiman2001random}, and XGBoost \citep{chen2016xgboost}, under
both a fixed training-set-size regime and a regime where training set size scales with
$D$.

Our contributions are:
\begin{itemize}
\item A synthetic benchmark generator with a controllable, low-dimensional nonlinear
      true signal embedded in up to 10,000 ambient features, released with this paper.
\item An empirical characterization of how held-out test RMSE/R\textsuperscript{2} for
      TabPFN v2 and TabICL scale with $D$, relative to tree and linear baselines, under
      both fixed-$N$ and $N$-scaling-with-$D$ regimes.
\item A direct empirical test of the feature-bagging ensemble workaround
      \citep{feuer2023scaling} at the specific scale (10,000 features) it was proposed
      to address.
\end{itemize}

\section{Related Work}
\label{sec:related}

\textbf{Tabular foundation models.} \citet{muller2022transformers} introduced Prior-Data
Fitted Networks (PFNs), transformers trained on synthetic data to approximate Bayesian
inference in a single forward pass, which \citet{hollmann2025tabpfn} scaled into TabPFN
v2, a model pretrained once and applied zero-shot to new tabular datasets via
in-context learning, achieving Nature-cover accuracy on datasets up to 10,000 rows and
500 features. \citet{qu2025tabicl} proposed TabICL, which restructures the in-context
learning architecture (a column-then-row transformer with a fixed-size latent
summary) specifically to scale to far larger row counts and a higher feature cap.
\citet{ma2024tabdpt} (TabDPT) and \citet{zeng2025tabflex} (TabFlex) pursue related
scaling directions -- retrieval-augmented real-data pretraining, and linear-attention
mechanisms, respectively -- motivated by the same observation we study directly: that
naively applying an in-context tabular learner outside its trained regime degrades
performance.

\textbf{Scaling TabPFN past its trained regime.} \citet{feuer2023scaling} is the paper
most directly relevant to our method: they show TabPFN's accuracy falls as feature
count grows past its trained regime, and propose closing the gap with sketching,
feature selection, and feature-bagging ensembles -- fitting independent copies of
TabPFN on random feature subsets and averaging. We adopt their feature-bagging
approach as the mechanism for applying TabPFN v2 and TabICL beyond their native caps,
and, unlike their work (which evaluates on existing high-dimensional tabular
benchmarks, where the number of truly informative features is unknown and grows with
$D$), we hold the number of truly informative features fixed at 5 while sweeping $D$,
isolating the effect of ambient dimensionality from the effect of additional signal.

\textbf{High-dimensional tabular learning more broadly.} Feature selection, random
projections, and ensembling over feature subsets are classical tools for
high-dimensional tabular learning with tree and linear models; our contribution is not
a new such tool, but an empirical measurement of how well the existing tool
(feature-bagging) closes the dimensionality gap for the specific architectures that
motivated it.

\section{Method}
\label{sec:method}

\subsection{Synthetic data generating process}
\label{sec:dgp}

Each dataset is parameterized by an ambient feature count $D \geq 5$ and a sample
count $N$. We draw $D_{\text{true}}=5$ latent features
$z = (z_1, \dots, z_5) \sim \mathcal{N}(0, I_5)$ per row and define the regression
target through a fixed nonlinear function:
\begin{align}
f(z) &= 3\sin(z_1 z_2) + 2 z_3^2 - 1.5\, z_4 z_5 + z_1 z_4^2 \nonumber \\
     &\quad - 0.5\, z_2^3 + 2\tanh(z_1 + z_3 - z_4), \\
y &= f(z) + \varepsilon, \qquad \varepsilon \sim \mathcal{N}(0, (\eta \cdot \sigma_{f})^2),
\end{align}
where $\sigma_f$ is the empirical standard deviation of $f(z)$ over the batch and
$\eta$ (the noise-to-signal ratio) is fixed across all $D$ so the achievable ceiling on
test $R^2$ does not itself vary with dimensionality.

The remaining $D-5$ ambient features are split into two kinds. A fraction
$p_{\text{noise}}$ are pure noise, $x_j \sim \mathcal{N}(0,1)$, independent of $y$. The
rest are noisy derived features: for each, we sample $k \in \{1,2,3\}$ of the true
latent features uniformly without replacement, a random weight vector, and one of four
transforms (identity, square, sine, tanh) applied to the weighted sum, then add
independent Gaussian noise at scale $\rho$. These derived features are correlated with
$y$ (through $z$) but not perfectly recoverable, matching the informal specification
that ambient features be "correlated or noisy functions of the true features...but not
100\% derivable." All $D$ columns (true, derived, and pure-noise) are then randomly
permuted so column position carries no information. Train, validation, and test splits
are drawn as one i.i.d. block per (D, N, seed) and sliced, since rows are independent
draws.

\subsection{Dimensionality and sample-size grid}
\label{sec:grid}

We sweep $D \in \{10, 20, 50, 100, 200, 500, 1000, 2000, 5000, 10000\}$. Two regimes
control the training set size $N_{\text{train}}$: \textbf{fixed-N}, where
$N_{\text{train}}$ is held at a constant TODO across all $D$; and \textbf{scaling-N},
where $N_{\text{train}} = \text{clip}(4D, \text{TODO}, \text{TODO})$. Validation and
test set sizes (TODO, TODO) are held fixed across all conditions so evaluation is
comparable. Each cell is repeated over TODO random seeds.

\subsection{Models}
\label{sec:models}

We evaluate Ridge regression \citep{hoerl1970ridge}, Random Forest
\citep{breiman2001random}, XGBoost \citep{chen2016xgboost}, TabPFN v2
\citep{hollmann2025tabpfn}, and TabICL \citep{qu2025tabicl}. For the two foundation
models, whenever $D$ exceeds the model's native feature cap ($c=500$ for TabPFN v2,
$c=2000$ for TabICL), we apply a feature-bagging ensemble \citep{feuer2023scaling}:
partition the $D$ columns into $\lceil D/c \rceil$ disjoint chunks of size $\leq c$ via
a random permutation, fit an independent copy of the model on each chunk, and average
the chunks' predictions. When $D \leq c$ this reduces to a single chunk containing
every feature, i.e.\ the model runs exactly as it natively would.

\section{Experimental Setup}
\label{sec:setup}

TODO: hardware/software (Apple Silicon CPU, no CUDA), library versions, and any
model-specific settings (e.g. \verb|ignore_pretraining_limits=True| for TabPFN so
a single chunk that exceeds its internal threshold does not raise rather than degrade).

\section{Results}
\label{sec:results}

TODO: figures (\verb|figures/r2_vs_dimensionality.pdf|,
\verb|figures/rmse_vs_dimensionality.pdf|) and summary tables, once the sweep has
run to completion.

\section{Discussion and Limitations}
\label{sec:discussion}

TODO. Points to cover regardless of the empirical outcome:
\begin{itemize}
\item The feature-bagging ensemble is an approximation we impose on TabPFN/TabICL, not
      a capability native to either model; results past each model's native cap
      characterize \emph{our wrapper}, not the released model in isolation.
\item All experiments ran on CPU (Apple Silicon, no CUDA); TabPFN v2/TabICL are
      primarily validated and optimized for GPU inference, so absolute wall-clock
      numbers here are not representative of production latency.
\item TabPFN v2 and later PriorLabs checkpoints are released under a non-commercial
      research license; TabICL is BSD-3-Clause.
\item A single fixed nonlinear true function and one specific derived-feature
      construction were used; results characterize this generative family, not
      dimensionality scaling in general.
\end{itemize}

\section{Conclusion}
\label{sec:conclusion}

TODO.

\section*{Acknowledgements}
TODO or remove.

\bibliography{refs}
\bibliographystyle{icml2025}

\end{document}
